\chapter{来源、材料尺度与争议}
\label{app:sources}

\section{先问材料记录了谁的行动}
\label{sources:index}

本卷以《明实录》与《明史》维持中枢编年骨架，以奏疏、会典、地方志、碑刻、契约、族谱、考古材料及朝鲜、越南、蒙古、日本、琉球和欧洲记录补足不同问题。材料不是可以叠加成全知视角的积木。诏令最适合说明朝廷决定发布什么；奏疏可以看见官员如何提出问题；地方文书偶尔能进入土地、债务、工程和家庭；外部记录能校正北京的分类。每一类材料也都有自己的沉默。

\section{怎样沿材料回到正文}

本附录中的每一条书目路线都按四个问题安排：它是什么版本或材料类型、最适合追问哪个时段或问题、它在本卷进入了哪些叙事，以及它不能独自裁定什么。读者可先由本页的材料锚点进入相应路线，再回到\hyperref[app:source-index]{来源索引}所列的首次叙事；不要把书目条目或事件账本当作另一部年表。

本卷只列入来源矩阵和事件账本已经使用的材料群或书名。电子全文可帮助定位，但具体判断仍应回到可核的点校本、影印本、档案号、碑刻著录或发掘报告；同名书的不同修本、抄本与现代整理本不能不加区别地互换。

\section{实录的时间与改写的时间}

《明实录》提供极密的日期、任命、诏令和奏报线索，却是由朝廷机构保存、修纂的记录。洪武、永乐和嘉靖的实录不能自动说明命令抵达每一个地方的时间；英宗实录成于复位后的天顺朝，阅读土木、景泰与夺门时尤其要留意胜者重排的合法性语言。崇祯朝没有存世实录，《崇祯长编》、后出的正史、奏疏、地方志与邻国材料必须并读，不能被伪装成一部完整的现场日记。

\section{数字不是没有来历的事实}

军数、斩获、灾民、田亩、饷额和贸易额，常在同一事件中出现多种数值。军报有邀功、请饷和问责的处境；清丈与赋役册有征收、比较和隐漏的用途；地方志又可能在修纂时把不同年代的材料合并。正文因此区分法定额、报解额、实解额与实际负担，也避免把一场战役的战报数字或一处灾区的损失推成全帝国的统计。

\section{海洋与边疆的多方叙事}

海禁、朝贡、互市和“倭寇”等词首先是行政分类。明方实录、朝鲜王朝实录、琉球文书、越南编年、蒙古汗统材料和葡萄牙记录，往往在地点、身份、交易与胜负上并不一致。它们的差异不必立刻消除：不同政权关心的对象本就不同。哈密、建州、交趾、朝鲜、东南港口和辽东的叙事，必须让多方材料各自说明其可证明的部分。

\section{地方社会如何只留下局部声音}

黄册、鱼鳞图册、契约、碑刻、宗谱、城墙遗址和寺观文献，让国家文书之外的生活偶尔可见；它们保存不均，且多出自有书写资源的家庭、机构和地方。妇女、雇工、佃户、军户、船户与流民常只以关系、税役或诉讼当事人的方式出现。承认这种缺口，不等于放弃社会史；它要求我们不以一县、一族或一份精英文集替全体人说话。

\section{基础书目与读法索引}
\label{sources:reading}

\subsection{编年、正史与晚明线索}

\begin{description}
  \item[\sourceanchor{ming-shilu}{《明实录》：太祖至熹宗各朝实录}] \textbf{版本/类型：}逐朝官修实录；本卷以朝别、年号、月日条定位，不将不同朝实录混作一部同时记录。 \textbf{适用：}1368--1627 年的诏令、任免、奏议摘要、军政和对外交涉。 \textbf{本卷用法：}为\eventref{EM-1427-JIAOZHI-WITHDRAWAL}{撤出交趾}、\eventref{EM-1449-TUMU}{土木堡}、海防、清丈和辽东等事件建立中央时间线。 \textbf{不能承担：}地方是否执行、战果和户口数字的准确性，或皇帝决策过程的全貌；建国、靖难、复辟与党争还须注意修纂时的选择。
  \item[\sourceanchor{ming-shi}{《明史》：本纪、志、表、列传}] \textbf{版本/类型：}清初官修正史的不同卷类；本卷按本纪、食货志、兵志、外国传及相关列传分别使用。 \textbf{适用：}1368--1644 年的制度线、人物结局、地理与崇祯年间的连续线索。 \textbf{本卷用法：}同实录互校，并为\eventref{ML-1644-902}{1644 年北京易手}及实录以外的制度和人物问题提供后出汇整。 \textbf{不能承担：}事件发生时的独立见证，也不能把清初忠奸、流寇、东林、宦官或亡国原因的判词直接当作解释。
  \item[谈迁《国榷》] \textbf{版本/类型：}明清之际的私家编年汇辑，本卷按可核年条与实录、奏疏互读。 \textbf{适用：}1368--1627 年，尤其是晚明宫廷与政争的异文、顺序和遗漏。 \textbf{本卷用法：}补充晚明事件链的检索路线。 \textbf{不能承担：}崇祯朝的材料替代品，或未经回查的日期、动机和数字。
  \item[\sourceanchor{chongzhen-changbian}{《崇祯长编》}] \textbf{版本/类型：}后出的崇祯朝编年辑录，须逐条辨明其所据材料。 \textbf{适用：}1628--1644 年的日序与朝局线索。 \textbf{本卷用法：}与奏疏、地方文献和对方记录一起阅读\eventref{ML-1629-901}{己巳之变}、\eventref{ML-1637-901}{练饷}和\eventref{ML-1644-902}{北京易手}。 \textbf{不能承担：}崇祯实录的地位；命令传达、军数和具体日序无法互证时仍保留分歧。
  \item[《明季北略》《明季南略》及晚明见闻、遗民笔记] \textbf{版本/类型：}鼎革前后形成或整理的见闻与记忆文本。 \textbf{适用：}约 1621--1644 年的京师、战乱、起义和易代记忆。 \textbf{本卷用法：}为晚明危机提供候选线索和版本比较。 \textbf{不能承担：}孤证精确军数、密谋、死亡数字或人物品格。
\end{description}

\subsection{制度、财政与官员文书}

\begin{description}
  \item[\sourceanchor{statutes}{《明会典》各修本}] \textbf{版本/类型：}永乐、弘治、嘉靖、万历等阶段的中央制度汇编，引用须分辨修纂时点和卷次。 \textbf{适用：}官制、仪制、赋役、机构职责和制度名目。 \textbf{本卷用法：}从制度文本进入\eventref{MM-1578-NATIONWIDE-SURVEY}{清丈}、\eventref{MM-1580-ONE-WHIP-ORDERS}{一条鞭}与军饷问题。 \textbf{不能承担：}某项制度在各地已经实行，或法定额就是实征额和实际负担。
  \item[《大诰》《大明律》、问刑条例与诏令汇编] \textbf{版本/类型：}法令、条例和诏令的规范性文本。 \textbf{适用：}洪武法制及其后的罪名、法定程序与制度变化。 \textbf{本卷用法：}辨认国家规定了什么，而不把规定写成社会现实。 \textbf{不能承担：}基层知法程度、判决率或社会行为；这些问题须另见案牍、地方志、碑刻或文集。
  \item[\sourceanchor{memorials}{奏疏、题本、揭帖、邸报与《内阁大库档案》所存明代档案}] \textbf{版本/类型：}有发文机关、收发和批答链的行政文书；引用应保留原件或抄件、档号及缺页状况。 \textbf{适用：}特别是保存较多的晚明财政、军政、司法与地方事务。 \textbf{本卷用法：}把题请、奉准、转饬和可见执行拆开阅读\eventref{ML-1637-901}{练饷}及其他政策。 \textbf{不能承担：}奏报数字是中立调查，或一纸奉准已证明地方完成执行。
  \item[《皇明经世文编》、官员文集与奏议集] \textbf{版本/类型：}经世选编及部院、官员的论政文本。 \textbf{适用：}明中后期财政、河工、海防、军政和灾荒的当时论证。 \textbf{本卷用法：}追问问题如何被提出，以及方案在何处遭遇资源与行政限制。 \textbf{不能承担：}建议已经实施，或作者所见地域可以代表全境。
  \item[《万历会计录》] \textbf{版本/类型：}会计与收支记录，须保留年度、机构、预算或决算口径及缺项。 \textbf{适用：}万历财政、军饷和加派的特定年度或机构问题。 \textbf{本卷用法：}为\eventref{MM-1578-NATIONWIDE-SURVEY}{清丈}之后的财政阅读提供口径检查，并与\eventref{ML-1637-901}{练饷}的征解问题对读。 \textbf{不能承担：}单凭一张表断定财政总崩溃，或推算民众的平均负担。
\end{description}

\subsection{地方、物质与文化材料}
\sourceanchor{local-records}{}

\begin{description}
  \item[府、州、县志及卫所、山川、寺观专志] \textbf{版本/类型：}地方志书；须辨别明修底本、清修增补、纂者和修纂年份。 \textbf{适用：}特定地点的建置、灾异、赋役、河工、学校、市场和兵灾。 \textbf{本卷用法：}同中央文书对读\eventref{MM-1556-SHAANXI-EARTHQUAKE}{华县地震}、\eventref{MM-1465-JINGXIANG-UPRISING}{荆襄}等地方后果。 \textbf{不能承担：}未留材料地区的经验，或由一府一县推出全国趋势。
  \item[碑刻、墓志、祠记、桥闸堤坝题记] \textbf{版本/类型：}原碑、拓本或可靠著录中的金石材料；需说明发现地、重刻与释文异同。 \textbf{适用：}地点、日期、官衔、营建、捐输、宗教和地方组织。 \textbf{本卷用法：}校验工程和地方网络的可见痕迹。 \textbf{不能承担：}纪功文字所称的工程效能，或精英自述以外群体的处境。
  \item[家谱、族谱、分家文书、契约、账簿与诉讼、赋役文书] \textbf{版本/类型：}家族保存和区域文书，尤其须标出原件形成期、续修期与保存链。 \textbf{适用：}家庭迁徙、婚姻、继承、土地、债务、租佃和基层行政。 \textbf{本卷用法：}把一条鞭、田赋、劳役和地方互助放回特定地区的实践，而非抽象的全国制度。 \textbf{不能承担：}一姓、一份契约或一个富庶地区代表全省或全国。
  \item[\sourceanchor{material-evidence}{城址、关隘、卫所、港口、沉船、窑址、钱币、银锭与墓葬考古}] \textbf{版本/类型：}有调查或发掘单位、层位、断代方法和正式报告的物证。 \textbf{适用：}营建、交通、技术、物资流动与边地聚落的物质条件。 \textbf{本卷用法：}为\eventref{ML-1606-901}{海龙屯}及边防、港口和生产问题检验空间可行性。 \textbf{不能承担：}单独复原诏令动机、具体谈判、战役路线或政治程序。
  \item[文集、笔记、书画题跋、戏曲小说、科举录与书院碑记] \textbf{版本/类型：}作者文本、选本、题跋和教育记录，须辨别成书、刊刻和重刻时间。 \textbf{适用：}士人交游、教育、出版、城市消费、知识传播与灾乱感受。 \textbf{本卷用法：}进入\eventref{MM-1579-ACADEMY-RESTRICTIONS}{书院限制}等文化和制度问题。 \textbf{不能承担：}文学虚构或规范性话语就是政治纪实、识字率或普遍思想。
\end{description}

\subsection{边疆、海洋与外部观察}
\sourceanchor{external-records}{}

\begin{description}
  \item[《筹海图编》及海防、倭患专书] \textbf{版本/类型：}十六、十七世纪官员和士人编纂的海防知识、方案与图说。 \textbf{适用：}东南海防、岛屿、船政、军器和倭患的地名与议题。 \textbf{本卷用法：}同地方与外部材料阅读\eventref{MM-1523-NINGBO-TRIBUTE-CONFLICT}{宁波朝贡冲突}、\eventref{MM-1547-ZHU-WAN-APPOINTMENT}{浙闽海防}和\eventref{MM-1567-YUEGANG-OPENING}{月港开海}。 \textbf{不能承担：}地图就是实测疆界、兵书就是实际编制，或“倭寇”是单一族群。
  \item[马欢《瀛涯胜览》、费信《星槎胜览》与郑和时代碑刻] \textbf{版本/类型：}随行者的航海见闻及可定位的碑刻。 \textbf{适用：}十五世纪印度洋航程、港口、使团和朝贡秩序。 \textbf{本卷用法：}为早明海外联系提供使团所见与官方秩序之间的对照。 \textbf{不能承担：}精确复原当地政治、贸易规模或所有航行参与者的经验。
  \item[《朝鲜王朝实录》与《历代宝案》] \textbf{版本/类型：}朝鲜王朝官修实录和琉球外交文书汇编；须保留各自历法、译名和抄传层次。 \textbf{适用：}朝贡、使节、海上贸易、倭患与跨海战争。 \textbf{本卷用法：}从非北京的观察位置对读\eventref{MM-1523-NINGBO-TRIBUTE-CONFLICT}{宁波}、\eventref{ML-1592-902}{援朝调兵}与东南港口。 \textbf{不能承担：}自动裁决两边记录不合的日期、身份、交易和胜负。
  \item[《大越史记全书》《蒙古源流》及清太祖、清太宗相关实录] \textbf{版本/类型：}越南、蒙古与后金／清方的编年、汗统或官修记录；须说明所据译本、历法和政治立场。 \textbf{适用：}交趾、北方草原、建州与辽东的跨政权行动。 \textbf{本卷用法：}校读\eventref{EM-1427-JIAOZHI-WITHDRAWAL}{撤出交趾}、\eventref{EM-1449-TUMU}{土木堡}、\eventref{ML-1618-901}{辽东攻势}和\eventref{ML-1621-901}{辽阳失守}。 \textbf{不能承担：}把任何一方的族名、疆界、兵数和胜败分类直接当成中性事实。
  \item[葡萄牙、西班牙、荷兰旅行记、书信、商馆与公司档案；利玛窦《中国札记》] \textbf{版本/类型：}海外旅行、传教、贸易与公司记录，含翻译、转述和殖民利益的过滤。 \textbf{适用：}十六、十七世纪华南、澳门、马尼拉、日本和海上贸易网络。 \textbf{本卷用法：}为\eventref{ML-1622-902}{澳门海防}与\eventref{ML-1624-901}{澎湖交涉}补入海外观察。 \textbf{不能承担：}中国地方社会的替代档案，或从外来印象直接推出贸易统计与社会常态。
\end{description}
